\section{Conclusions}\label{sec:conclusions}

We presented a complex128 statevector engine for one, two and four AMD Alveo U55C cards and measured it, together with
a baseline design that computes the gates on partition qubits on the host, on four physical cards in one host. The engine executes every gate that needs no
amplitudes of another slice on the cards, including diagonal, slice controlled and slice permutation gates on partition
qubits and multi controlled gates wider than its tile, packs consecutive gates into HBM passes, and spreads the
state over all 32 pseudo channels of every card. It simulates \NumMaxWidthOneCard~qubits on one card, \NumMaxWidthTwoCards~on two and \NumMaxWidthFourCards~on four cards,
every instance agreeing with an independent reference to double precision rounding (largest amplitude error
\NumWorstMaxAbsErr). On the same cards the QFT at $n=20$ runs \NumSubmittedOverRevisedQftTwenty~times faster than in the
baseline design, whose timing curve a microbenchmark model without fitted parameters reproduces to within a factor
of \NumModelRatioMin~to \NumModelRatioMax.

The measurements also set the limits. Qiskit Aer on the \NumHostCores~core host is faster than the four card system for
almost all circuits measured. More cards add capacity but not speed, because remaps that involve a card qubit pass
through host memory: in \NumScaleBelowOne~of \NumScaleCases~comparisons two or four cards are slower than one, and the
energy of a simulation grows with the number of cards. One card of the host limits the widest runs thermally. The next
steps are therefore a faster exchange path between the cards, by peer to peer transfers or direct links, and a kernel
with a higher rate per card.
